% GENERATED FILE. Edit canonical lessons, then run npm run generate:books.
% canonical-source-hash: fnv1a64:2a1f8919af88fcfb
% canonical-lessons: LA-C342-fugo, LA-C342-supero, LA-C342-caedo, LA-C342-vasto, LA-C342-clades

\chapter{Clādēs}
\label{ch:la-clades}
\hlchaptermodality{342}

\begin{chapteropening}
\textbf{By the end of this chapter:} \emph{I can tell how a battle ended: to put to flight, to overcome, to cut down, to lay waste and a disaster.}

Complete the last lesson of chapter 342: I can tell how a battle ended: to put to flight, to overcome, to cut down, to lay waste and a disaster.
\end{chapteropening}

\section[fugō, fugāre]{fugō, fugāre — to put to flight, to rout}

\label{lesson:LA-C342-fugo}



\begin{quote}
Before the new one: say the Latin for to take by storm, then the Latin for a battering ram.
\end{quote}

\subsection*{You'll want to know: fugō, fugāre}
\textbf{fugō, fugāre} — \textquotedblleft{}to put to flight, to rout\textquotedblright{}.

\begin{quote}
Equitēs nostrī hostēs fugāvērunt.
\end{quote}

\begin{quote}
\emph{Our cavalry put the enemy to flight.}
\end{quote}

\emph{fugō, fugāre, fugāvī, fugātum}: to make others run.

\begin{cousinweb}
\textbf{fugō} is built on \emph{fuga}, \textquotedblleft{}flight\textquotedblright{} (Lewis \& Short): to cause to flee.
\end{cousinweb}

\begin{grammarlens}[title={What you've built}]
One more word for telling, explaining and describing at B1.
\end{grammarlens}

\subsection*{Your turn}
\begin{itemize}
\raggedright
  \item \emph{Say it:} \emph{fugō, fugāre}
  \item \emph{Say it:} \emph{fugō, fugāre}, once more
  \item \emph{Say it:} a sentence of your own with \emph{fugō, fugāre}
  \item \emph{Recall:} say the Latin for to take by storm, then the Latin for a battering ram, then say \emph{fugō, fugāre} again
\end{itemize}

\begin{tcolorbox}[breakable,colback=teal!4,colframe=teal!35!black,title={Before you move on}]
What does fugō, fugāre mean? (\textquotedblleft{}To put to flight, to rout\textquotedblright{}.) Say it once more.
\end{tcolorbox}
\section[superō, superāre]{superō, superāre — to overcome, to defeat; to surpass}

\label{lesson:LA-C342-supero}



\begin{quote}
Before the new one: say the Latin for a battering ram, then the Latin for to put to flight.
\end{quote}

\subsection*{You'll want to know: superō, superāre}
\textbf{superō, superāre} — \textquotedblleft{}to overcome, to defeat; to surpass\textquotedblright{}.

\begin{quote}
Parva manus magnum exercitum superāvit.
\end{quote}

\begin{quote}
\emph{A small band defeated a large army.}
\end{quote}

\emph{superō, superāre, superāvī, superātum}, from \emph{super}, 'above'. Also to climb over: \emph{montēs superāre}.

\begin{cousinweb}
\textbf{superō} is built on \emph{super}, \textquotedblleft{}above, over\textquotedblright{} (Lewis \& Short): to go over, to get the upper hand.
\end{cousinweb}

\begin{grammarlens}[title={What you've built}]
One more word for telling, explaining and describing at B1.
\end{grammarlens}

\subsection*{Your turn}
\begin{itemize}
\raggedright
  \item \emph{Say it:} \emph{superō, superāre}
  \item \emph{Say it:} \emph{superō, superāre}, once more
  \item \emph{Say it:} a sentence of your own with \emph{superō, superāre}
  \item \emph{Recall:} say the Latin for a battering ram, then the Latin for to put to flight, then say \emph{superō, superāre} again
\end{itemize}

\begin{tcolorbox}[breakable,colback=teal!4,colframe=teal!35!black,title={Before you move on}]
What does superō, superāre mean? (\textquotedblleft{}To overcome, to defeat; to surpass\textquotedblright{}.) Say it once more.
\end{tcolorbox}
\section[caedō, caedere]{caedō, caedere — to cut down, to kill; to beat}

\label{lesson:LA-C342-caedo}



\begin{quote}
Before the new one: say the Latin for to put to flight, then the Latin for to overcome.
\end{quote}

\subsection*{You'll want to know: caedō, caedere}
\textbf{caedō, caedere} — \textquotedblleft{}to cut down, to kill; to beat\textquotedblright{}.

\begin{quote}
Mīlitēs hostēs fugientēs caedēbant.
\end{quote}

\begin{quote}
\emph{The soldiers were cutting down the fleeing enemy.}
\end{quote}

\emph{caedō, caedere, cecīdī, caesum}. Also to fell trees and to beat. Its noun is \emph{caedēs}.

\begin{cousinweb}
\textbf{caedō} is inherited from a Proto-Indo-European root meaning \textquotedblleft{}to cut, to hew\textquotedblright{} (EtymDB). English words in \emph{-cide}, such as \emph{homicide}, come from its family.
\end{cousinweb}

\begin{grammarlens}[title={What you've built}]
One more word for telling, explaining and describing at B1.
\end{grammarlens}

\subsection*{Your turn}
\begin{itemize}
\raggedright
  \item \emph{Say it:} \emph{caedō, caedere}
  \item \emph{Say it:} \emph{caedō, caedere}, once more
  \item \emph{Say it:} a sentence of your own with \emph{caedō, caedere}
  \item \emph{Recall:} say the Latin for to put to flight, then the Latin for to overcome, then say \emph{caedō, caedere} again
\end{itemize}

\begin{tcolorbox}[breakable,colback=teal!4,colframe=teal!35!black,title={Before you move on}]
What does caedō, caedere mean? (\textquotedblleft{}To cut down, to kill; to beat\textquotedblright{}.) Say it once more.
\end{tcolorbox}
\section[vāstō, vāstāre]{vāstō, vāstāre — to lay waste, to ravage}

\label{lesson:LA-C342-vasto}



\begin{quote}
Before the new one: say the Latin for to overcome, then the Latin for to cut down.
\end{quote}

\subsection*{You'll want to know: vāstō, vāstāre}
\textbf{vāstō, vāstāre} — \textquotedblleft{}to lay waste, to ravage\textquotedblright{}.

\begin{quote}
Hostēs agrōs vīllāsque vāstāvērunt.
\end{quote}

\begin{quote}
\emph{The enemy laid waste the fields and the farmhouses.}
\end{quote}

\emph{vāstō, vāstāre, vāstāvī, vāstātum}: to make a land empty. English \emph{devastate}.

\begin{cousinweb}
\textbf{vāstō} is built on \emph{vāstus}, \textquotedblleft{}empty, desolate\textquotedblright{} (EtymDB, Lewis \& Short): to make a place empty. English \emph{waste} and \emph{devastate} come from its family.
\end{cousinweb}

\begin{grammarlens}[title={What you've built}]
One more word for telling, explaining and describing at B1.
\end{grammarlens}

\subsection*{Your turn}
\begin{itemize}
\raggedright
  \item \emph{Say it:} \emph{vāstō, vāstāre}
  \item \emph{Say it:} \emph{vāstō, vāstāre}, once more
  \item \emph{Say it:} a sentence of your own with \emph{vāstō, vāstāre}
  \item \emph{Recall:} say the Latin for to overcome, then the Latin for to cut down, then say \emph{vāstō, vāstāre} again
\end{itemize}

\begin{tcolorbox}[breakable,colback=teal!4,colframe=teal!35!black,title={Before you move on}]
What does vāstō, vāstāre mean? (\textquotedblleft{}To lay waste, to ravage\textquotedblright{}.) Say it once more.
\end{tcolorbox}
\section[clādēs, clādis]{clādēs, clādis — a defeat, a disaster}

\label{lesson:LA-C342-clades}



\begin{quote}
Before the new one: say the Latin for to cut down, then the Latin for to lay waste.
\end{quote}

\subsection*{You'll want to know: clādēs, clādis}
\textbf{clādēs, clādis} — \textquotedblleft{}a defeat, a disaster\textquotedblright{}.

\begin{quote}
Post clādem paucī mīlitēs in castra rediērunt.
\end{quote}

\begin{quote}
\emph{After the defeat few soldiers returned to the camp.}
\end{quote}

Feminine, 3rd declension. A heavy defeat with great loss of life.

\begin{cousinweb}
\textbf{clādēs} is usually traced to a root meaning \textquotedblleft{}to beat, to break\textquotedblright{}: EtymDB gives a Proto-Indo-European root of that sense, and Lewis \& Short compare Greek \emph{klaō}, \textquotedblleft{}to break\textquotedblright{}. The details are not settled.
\end{cousinweb}

\begin{grammarlens}[title={What you've built}]
One more word for telling, explaining and describing at B1.
\end{grammarlens}

\subsection*{Your turn}
\begin{itemize}
\raggedright
  \item \emph{Say it:} \emph{clādēs, clādis}
  \item \emph{Say it:} \emph{clādēs, clādis}, once more
  \item \emph{Say it:} a sentence of your own with \emph{clādēs, clādis}
  \item \emph{Recall:} say the Latin for to cut down, then the Latin for to lay waste, then say \emph{clādēs, clādis} again
\end{itemize}

\begin{tcolorbox}[breakable,colback=teal!4,colframe=teal!35!black,title={Before you move on}]
What does clādēs, clādis mean? (\textquotedblleft{}A defeat, a disaster\textquotedblright{}.) Say it once more.
\end{tcolorbox}
